% HANDOFF_TRACK_B.tex
% Self-contained document for the second student, describing the code+experiments
% work remaining on the TMLR revision (Track B).
% Compiles with pdflatex: pdflatex HANDOFF_TRACK_B; pdflatex HANDOFF_TRACK_B
\documentclass[11pt,a4paper]{article}

\usepackage[margin=2.5cm]{geometry}
\usepackage[T1]{fontenc}
\usepackage{lmodern}
\usepackage{microtype}
\usepackage{amsmath}
\usepackage{amssymb}
\usepackage{hyperref}
\usepackage{booktabs}
\usepackage{longtable}
\usepackage{array}
\usepackage{enumitem}
\usepackage{xcolor}
\usepackage{listings}

\hypersetup{
  colorlinks=true,
  linkcolor=blue!60!black,
  urlcolor=blue!60!black,
  citecolor=blue!60!black
}

\lstset{
  basicstyle=\ttfamily\small,
  breaklines=true,
  columns=fullflexible,
  backgroundcolor=\color{gray!8},
  frame=single,
  framerule=0pt,
  xleftmargin=6pt,
  xrightmargin=6pt,
  aboveskip=8pt,
  belowskip=8pt
}

\newcommand{\code}[1]{\texttt{#1}}
\newcommand{\file}[1]{\texttt{#1}}

\title{Track B --- Handoff Document\\[4pt]
       \large Experiment-Running Student\\[2pt]
       \normalsize TMLR Revision, TraLO Paper}
\author{Prepared by the advisor + Claude (writing track)}
\date{2026-07-21}

\begin{document}
\maketitle

\section*{Quick context}

The TMLR paper on TraLO went through a multi-perspective simulated review on
2026-07-21. Decision: \textbf{Major Revision}. \textbf{Track A} (edits to
the paper text) has been completed by the advisor with Claude. \textbf{Track
B} (new experiments) --- this document --- is handed off to you.

\begin{itemize}[nosep]
  \item \textbf{Active paper source}: \file{tmlr\_submission/main.tex} in
        this project, plus the Overleaf working copy the advisor will invite
        you to.
  \item \textbf{Full revision roadmap}: \file{PLAN.md} in the project root.
  \item \textbf{The 1944-run corpus}: already exists at the location agreed
        with the advisor (not inside the \file{tmlr\_submission/} directory).
  \item \textbf{Acceptance probability status}:
    \begin{itemize}[nosep]
      \item After Track A (current): $\sim$62--64\%
      \item + B1 (imbalanced baselines): 70--75\%
      \item + B1 + B2 (native-resolution HAM10000): 78--85\%
      \item + full Track B: 82--88\%
    \end{itemize}
\end{itemize}

\section*{Prerequisites}

\begin{itemize}[nosep]
  \item Access to the training code (TraLO + Fioretto-LDF + Hounie-RCL +
        Shifman-LP).
  \item BIU SLURM: partition \code{H200-12h}, conda env \code{crs}.
        \textbf{Important}: do not run anything on \code{slurm-login1} ---
        only \code{sbatch} or connect to \code{dsisco01} / \code{dsisco02}
        via \code{ssh dsisco01} / \code{ssh dsisco02}.
  \item Overleaf: the advisor will invite you to the project.
  \item Python: for the analysis scripts (\code{scripts/} in the original
        code directory).
\end{itemize}

\section*{Work division while you run experiments}

\begin{itemize}[nosep]
  \item \textbf{You}: any code changes, SLURM runs, generation of new tables
        from results, updates to \code{main.tex} and \code{references.bib}
        in Overleaf reflecting new findings.
  \item \textbf{Advisor + Claude (local)}: if wording tweaks unrelated to
        results are needed during your work, they happen in the local files
        and are synced to Overleaf.
  \item \textbf{Sync}: any significant change --- commit in Overleaf. Try to
        update \file{PLAN.md} with a checkbox for the item you completed
        (or notify the advisor to do so).
\end{itemize}

% ============================================================================
\section{Priority-ordered Track B items}

\subsection*{Must-do --- without these, acceptance probability will not
             pass 65\%}

% ----------------------------------------------------------------------------
\subsection{B1. Imbalanced-learning baselines (highest priority)}

\paragraph{What is missing.} All five reviewers requested a baseline of
focal loss + clipping / class-balanced + clipping / logit adjustment +
clipping. Section 6 (Limitations) of the paper already concedes this gap.
The Devil's Advocate went further: if focal + clipping closes the
\mbox{+1.6--5.3~pp} macro-F1 gap over vanilla clipping, the entire macro-F1
claim collapses.

\paragraph{Goal.} Answer definitively: does an imbalanced-learning-aware
clipper close the macro-F1 gap the paper attributes to constraint-time
training?

\paragraph{Minimal scope.}
\begin{itemize}[nosep]
  \item 3 methods: focal loss (Lin et al.\ 2017), class-balanced weights
        (Cui et al.\ 2019), logit adjustment (Menon et al.\ 2020).
  \item 3 datasets: TissueMNIST, DermMNIST, OctMNIST.
  \item 3 backbones: MobileNetV3, RegNetY-400MF, ViT-B/16.
  \item 1 tight cap: L30 only (where the macro-F1 gap of TraLO is largest).
  \item 4 seeds (1--4, as in the rest of the paper).
  \item Post-hoc clipping applied via the same method as Shifman-LP.
  \item \textbf{Total}: $3 \times 3 \times 3 \times 4 = \mathbf{108}$
        \textbf{runs}.
\end{itemize}

\paragraph{Code location.}
\begin{itemize}[nosep]
  \item Add new methods in the same \code{struct} as the existing methods.
  \item Focal loss: $\ell = -\alpha (1-p)^\gamma \log(p)$, with
        $\alpha{=}0.25$, $\gamma{=}2.0$ (defaults from the original paper).
  \item Class-balanced: reweight class $c$ on-the-fly with
        $(1-\beta)/(1-\beta^{n_c})$, $\beta{=}0.9999$.
  \item Logit adjustment: add $\tau \log(\text{prior}_c)$ to logits,
        $\tau{=}1.0$.
\end{itemize}

\paragraph{Deliverables.}
\begin{enumerate}[nosep]
  \item Results CSV: \file{results/imbalanced\_baselines\_L30.csv} \\
        (columns: \code{dataset, backbone, seed, method, macro\_f1, cc\_f1,
        native\_satisfied, flips}).
  \item New table: \file{tables/tab\_imbalanced\_baselines.tex} \\
        Suggested structure: rows = dataset$+$backbone, columns = 3
        imbalanced methods + TraLO + gap. The advisor + Claude can finalize
        the caption.
  \item Updates to \code{main.tex}:
    \begin{itemize}[nosep]
      \item Section~5 (Results): add a new subsection
            \code{\textbackslash{}subsection\{Imbalanced-training baselines
            with clipping\}} --- 1--2 paragraphs after Sec.~5.3 (deploy).
      \item Abstract: if TraLO still leads, update to ``over both vanilla
            clipping and imbalanced-training baselines''. If not, update to
            ``comparable overall quality to imbalanced-training baselines
            while additionally satisfying the cap natively''.
      \item Section~6 (Limitations): remove or shorten the paragraph
            \textit{Missing baselines from imbalanced learning}.
    \end{itemize}
\end{enumerate}

\paragraph{Estimated cost.} $\sim$9 GPU-hours + $\sim$half-day of coding
(adding 3 methods).

% ----------------------------------------------------------------------------
\subsection{B2. Native-resolution medical benchmark (HAM10000)}

\paragraph{What is missing.} Devil's Advocate CRITICAL --- all experiments
are on 28$\times$28 upsampled to 224$\times$224. There is no evidence that
the OctMNIST tight-cap phenomenon occurs at real resolution. If it does
not $\to$ the deployment framing in the paper drops to ``toy scale''.

\paragraph{Goal.} Show that the general behavior of TraLO in the tight-cap
regime reproduces at native resolution.

\paragraph{Minimal scope.}
\begin{itemize}[nosep]
  \item Dataset: HAM10000 at native resolution ($600\times450$) --- the
        source of DermMNIST.
  \item Constrained class: melanoma.
  \item 2 tight caps: L30, L40.
  \item 3 backbones: same as the main grid.
  \item 4 seeds.
  \item \textbf{Methods}: TraLO + Fioretto-LDF + Hounie-RCL + Heuristic
        clipper (4 methods total).
  \item \textbf{Total}: $4 \times 2 \times 3 \times 4 = \mathbf{96}$
        \textbf{runs} (or 24 for TraLO alone + 72 for baselines).
\end{itemize}

\paragraph{Technical notes.}
\begin{itemize}[nosep]
  \item Standard ImageNet preprocessing (256$\times$256 crop $\to$
        224$\times$224 center crop) is not directly relevant here --- do a
        resize to 450$\times$450 or the size suggested per backbone (see
        \code{torchvision.models}).
  \item Warmup: 50 epochs (same budget).
  \item Constraint phase: 300 epochs.
  \item Runtime: likely slower because of the larger image --- expect
        15--30 minutes per run.
\end{itemize}

\paragraph{Deliverables.}
\begin{enumerate}[nosep]
  \item \file{results/hamres\_native\_L30\_L40.csv}
  \item \file{tables/tab\_hamres\_native.tex}
  \item \file{figures/fig\_hamres\_octanalog.pdf} --- a plot analogous to
        \file{fig\_octmnist} (top: absolute values; bottom: paired delta).
  \item Updates to \code{main.tex}:
    \begin{itemize}[nosep]
      \item Section~5: new subsection
            \code{\textbackslash{}subsection\{Replication on
            native-resolution HAM10000\}} (or as an appendix).
      \item Abstract: if the phenomenon reproduces, remove ``on
            MedMNIST-scale tasks'' from the practitioners paragraph and
            replace with ``across two scales''.
      \item Section~6 (Limitations): shorten \textit{Dataset scope}
            accordingly.
    \end{itemize}
\end{enumerate}

\paragraph{Failure scenario.} If the phenomenon does \emph{not} reproduce
at native resolution, that is itself an important finding, added as
``native-resolution replication'' with the negative result reported.
Devil's Advocate would still be satisfied because the CRITICAL items were
about the absence of an experiment, not about a positive outcome.

\paragraph{Estimated cost.} $\sim$4--6 GPU-hours + one day of work on the
native HAM10000 loader.

\subsection*{Strongly recommended --- lifts from 75\% to 85\%}

% ----------------------------------------------------------------------------
\subsection{B3. ALM (Augmented Lagrangian) as a baseline}

\paragraph{Why.} R2 (Domain reviewer) pointed out that ALM is the standard
literature fix for linear-penalty windup. Section~2 of the paper now
justifies (after A12) why we did not run ALM --- but running ALM as a
baseline on 6 OctMNIST tight cells would close the concern definitively.

\paragraph{Scope.}
\begin{itemize}[nosep]
  \item 1 method: Fioretto-LDF with the ALM update rule:
        $\lambda_c \leftarrow \max(0, \lambda_c + \eta(S_c - K_c)) +
        \mu(S_c - K_c)^+$, with $\mu$ growing linearly.
  \item 6 cells: OctMNIST L30 + L40 $\times$ 3 backbones $\times$ 4 seeds
        $= \mathbf{24}$ \textbf{runs}.
\end{itemize}

\paragraph{Deliverables.}
\begin{itemize}[nosep]
  \item A new row in the graft table (or a separate table):
        \file{tables/tab\_alm.tex}.
  \item Update in Section~5 or App.~C.1 (graft) --- a short paragraph:
        ``ALM as a fifth baseline: results here.''
  \item If ALM does not close the gap, the A12 justification in Section~2
        stands. If it does, refine the A12 paragraph accordingly.
\end{itemize}

\paragraph{Estimated cost.} $\sim$2 GPU-hours + $\sim$half-day of coding.

% ----------------------------------------------------------------------------
\subsection{B4. Extend the component ablation to a tie-region cell}

\paragraph{Why.} R1 M4 + Devil's Advocate --- the ablation was run only on
tight-cap OctMNIST cells (the win region). Open question: are reset+hinge
load-bearing in tie regions too, or only at tight caps?

\paragraph{Scope.}
\begin{itemize}[nosep]
  \item One tie-region cell: DermMNIST L50 symmetric with MobileNetV3.
  \item Full ablation of 4 components (reset / hinge / $\rho$ schedule /
        lambda freeze).
  \item 4 seeds.
  \item \textbf{Total}: $4 \times 4 = \mathbf{16{-}24}$ \textbf{runs}
        (depending on how many variants).
\end{itemize}

\paragraph{Deliverables.}
\begin{itemize}[nosep]
  \item Additional rows in \file{tab\_ablation\_complete.tex} (or a
        separate table \file{tab\_ablation\_tie\_cell.tex}).
  \item Update in Section~5 (What Carries the Win): a short paragraph
        such as ``In a tie-region cell (DermMNIST L50), the reset and
        hinge components \dots{} [continue with actual result]''.
  \item One of two possible conclusions:
    \begin{itemize}[nosep]
      \item reset+hinge load-bearing in tie regions too $\to$ strengthens
            the claim that portable components carry the effect everywhere.
      \item reset+hinge not load-bearing in tie regions $\to$ tighter claim:
            ``portable components carry the tight-cap advantage
            specifically''.
    \end{itemize}
  \item \textbf{Both outcomes are good for the paper.}
\end{itemize}

\paragraph{Estimated cost.} $\sim$2 GPU-hours.

\subsection*{Analyses on the existing corpus --- no new runs}

% ----------------------------------------------------------------------------
\subsection{B5. Sensitivity sweep of the win-bar rule}

\paragraph{Goal.} Address R1 M1 --- does regime classification change under
different win thresholds?

\paragraph{What to do.}
\begin{itemize}[nosep]
  \item On the existing corpus (1944 runs + ablation/graft), re-run the
        ``win regime'' rule: mean paired cc-F1 gap $\ge \tau$
        \textsc{and} at least half of cells clear $\tau$, for
        $\tau \in \{+0.003, +0.005, +0.010\}$.
  \item Produce a table: rows = regime (Tissue-tight, Tissue-mid,
        Tissue-loose, Derm-tight, \dots); columns = $\tau_{003}$ /
        $\tau_{005}$ / $\tau_{010}$ win/tie/loss.
  \item Expected pattern: the OctMNIST tight-cap classification as a win
        should be stable at all thresholds; most other regions tie at all
        thresholds.
\end{itemize}

\paragraph{Deliverable.}
\begin{itemize}[nosep]
  \item \file{tables/tab\_winbar\_sensitivity.tex}
  \item Add to Section~4.2 (or App.~B): 2--3 sentences referencing the
        table.
  \item The paragraph added in A5 already claims that ``OctMNIST tight-cap
        survives generously wider thresholds'' --- this table will
        numerically confirm that.
\end{itemize}

\paragraph{Estimated cost.} 1--2 hours of Python scripting.

% ----------------------------------------------------------------------------
\subsection{B6. Bootstrap CIs for the headline cells}

\paragraph{Goal.} R1 M3 + Devil's Advocate --- $n=4$ seeds means the
ViT-B/16 L30 effect of $+0.081$ is only $\sim$1.8$\sigma$. Bootstrap CIs
would strengthen the statistical support for the main-headline cells.

\paragraph{What to do.}
\begin{itemize}[nosep]
  \item On the results of the 6 OctMNIST tight-cap cells $\times$ cc-F1.
  \item Bootstrap 10{,}000 resamples per cell (with replacement); compute
        95\% CI on cell means.
  \item Produce a short table or an additional row in
        \file{tab\_oct\_backbone.tex}.
\end{itemize}

\paragraph{Deliverable.}
\begin{itemize}[nosep]
  \item If CIs do not cross 0 $\to$ strengthens the win claim (especially
        for ViT).
  \item If CIs do cross 0 $\to$ this must be reported in a visible place;
        it changes the strength of the claim.
\end{itemize}

\paragraph{Estimated cost.} 30 minutes of scripting.

% ----------------------------------------------------------------------------
\subsection{B7. Formal BH-FDR correction}

\paragraph{Goal.} The A5 paragraph (Track A) already claims that the
results survive BH-FDR. B7 will compute the $q$-values explicitly for the
full family of comparisons.

\paragraph{What to do.}
\begin{itemize}[nosep]
  \item Collect all $p$-values reported in the paper (paired Wilcoxon,
        sign tests).
  \item Run BH-FDR with $q=0.05$ on several families:
    \begin{itemize}[nosep]
      \item[(a)] headline OctMNIST 6 cells $\times$ 2 components = 12
                 tests;
      \item[(b)] full symmetric grid 27 cells $\times$ 2 metrics = 54
                 tests;
      \item[(c)] ablation-graft 24 comparisons.
    \end{itemize}
  \item Produce a table: rows = comparison, columns = raw $p$,
        BH-adjusted $q$, verdict (significant at $q=0.05$?).
\end{itemize}

\paragraph{Deliverable.}
\begin{itemize}[nosep]
  \item Short table in the appendix or body.
  \item Extension of the A5 paragraph with the formal numbers.
\end{itemize}

\paragraph{Estimated cost.} 1 hour.

% ----------------------------------------------------------------------------
\subsection{B8. Fourth backbone for the ``backbone-general'' cc-F1 claim}

\paragraph{What is missing.} MobileNetV2 is already in the supplementary
grid for macro-F1 (\file{tab\_backbone\_generality.tex}). Extend it to
cc-F1 at OctMNIST tight caps as well.

\paragraph{Scope.}
\begin{itemize}[nosep]
  \item MobileNetV2 $\times$ OctMNIST L30 + L40 $\times$ 4 seeds $\times$
        TraLO + best trained baseline $= \mathbf{16}$ \textbf{runs}.
\end{itemize}

\paragraph{Deliverable.}
\begin{itemize}[nosep]
  \item Additional row in \file{tab\_oct\_backbone.tex}.
  \item Update the ``backbone-general'' phrasing in the body to ``across
        four backbones tested (including MobileNetV2)''.
\end{itemize}

\paragraph{Estimated cost.} 40 minutes of runtime.

% ============================================================================
\section{Track B cost summary}

\begin{table}[h]
\centering
\begin{tabular}{lrrl}
\toprule
Item & Runs & GPU-hours (@5 min) & Code effort \\
\midrule
B1 imbalanced baselines & 108 & $\sim$9h & Medium (3 methods) \\
B2 native-res HAM10000 & 96 & $\sim$4--6h & Medium (loader) \\
B3 ALM & 24 & $\sim$2h & Low (Fioretto variant) \\
B4 tie-region ablation & 24 & $\sim$2h & Low \\
B5 win-bar sensitivity & 0 & 0 & Python script \\
B6 bootstrap CIs & 0 & 0 & Python script \\
B7 BH-FDR computation & 0 & 0 & Python script \\
B8 MobileNetV2 tight-cap & 16 & $\sim$1h & Low \\
\midrule
\textbf{Total} & \textbf{$\sim$270} & \textbf{$\sim$18--20 h} & \\
\bottomrule
\end{tabular}
\end{table}

For reference: the existing grid is 1944 runs. Track B adds $\sim$14\%.

% ============================================================================
\section{Overleaf working conventions}

\begin{enumerate}[nosep]
  \item \textbf{Before starting an item}: pull / sync in Overleaf to
        receive the latest Track A changes from the advisor.
  \item \textbf{While running experiments}: do not update \code{main.tex}
        yet --- wait for results.
  \item \textbf{After results arrive}: create/update a table file in
        \file{tmlr\_submission/tables/}, then update \code{main.tex} with
        the relevant sentences.
  \item \textbf{Sync}: after each completed item, commit in Overleaf with
        a clear message (e.g.\ ``B1 imbalanced baselines: added
        tab\_imbalanced\_baselines.tex + Sec 5.4'').
  \item \textbf{If wording or interpretation is unclear}: open a
        conversation with the advisor. Do not guess what the paper should
        say --- the tone of the paper is delicate (honest concessions in
        specific places) and the advisor will know the right phrasing.
\end{enumerate}

% ============================================================================
\section{Points where you should consult the advisor}

\begin{itemize}[nosep]
  \item If a result closes a gap or refutes a claim $\to$ the abstract
        needs rephrasing.
  \item If an unexpected result appears (e.g., focal loss + clipping
        closes the TraLO gap) $\to$ this is a substantive change to the
        paper's framing.
  \item If B2 (native-res HAM10000) shows the phenomenon does not
        reproduce $\to$ decide whether to report it as a negative
        replication or to remove the deployment framing from the whole
        paper.
\end{itemize}

% ============================================================================
\section{Final file structure (state after Track B completion)}

\begin{lstlisting}
tmlr_submission/
|-- main.tex                        # + 2-4 new subsections
|-- references.bib                   # updated
|-- tables/
|   |-- tab_ccf1.tex                 # (Track A)
|   |-- tab_ablation_complete.tex    # + rows from B4
|   |-- tab_graft.tex
|   |-- tab_deploy.tex
|   |-- tab_backbone_generality.tex  # + MobileNetV2 row (B8)
|   |-- tab_oct_backbone.tex         # + bootstrap CIs (B6)
|   |-- tab_imbalanced_baselines.tex # <- new (B1)
|   |-- tab_hamres_native.tex        # <- new (B2)
|   |-- tab_alm.tex                  # <- new (B3)
|   |-- tab_winbar_sensitivity.tex   # <- new (B5)
|   `-- (rest as-is)
`-- figures/
    `-- fig_hamres_octanalog.pdf     # <- new (B2)
\end{lstlisting}

% ============================================================================
\section*{Sign-off}

When you finish each item, mark \checkmark in \file{PLAN.md} under
``Track B''. If you perform an item outside this list or feel the scope
has changed, update \file{PLAN.md} accordingly.

Good luck.

\end{document}
